% ============================================================
% Floodcaster: Demo Track Submission for SIGSPATIAL 2026
% 4 pages including references; demo track deadline June 28, 2026
%
% Cross-references the Applications Track paper (FloodRSCT) for
% method details; this paper focuses entirely on the mobile artifact
% and the live demonstration.
% ============================================================

\documentclass[sigconf,nonacm=true,review,authoryear]{acmart}

\settopmatter{printacmref=false}
\renewcommand\footnotetextcopyrightpermission[1]{}
\pagestyle{plain}

\setcitestyle{authoryear,round}

\usepackage[utf8]{inputenc}
\usepackage[T1]{fontenc}
\usepackage{amsmath}
\usepackage{booktabs}
\usepackage{graphicx}
\usepackage{xcolor}
\usepackage{url}

% Safety net for tight two-column lines: allow modest extra stretch
% before declaring overfull. Less aggressive than \sloppy.
\setlength{\emergencystretch}{1em}

\acmConference[SIGSPATIAL '26]{The 34th ACM SIGSPATIAL International
  Conference on Advances in Geographic Information Systems}{November
  3--6, 2026}{Riverside, CA, USA}
\acmYear{2026}
\copyrightyear{2026}

% Per SIGSPATIAL 2026 Demo Track CFP, [Demo] suffix is required at
% submission and removed at camera-ready.
\title{Floodcaster: A Mobile Operator System for Auditable
  Flood Crisis Decision Support [Demo]}

\author{Rudolph A. Martin}
\affiliation{%
  \institution{Next Shift Consulting LLC}
  \city{Madison}
  \state{Wisconsin}
  \country{USA}
}
\email{rudy@nextshiftconsulting.com}

\begin{document}

% ============================================================
% ABSTRACT (≈100 words; demo abstracts are short)
% ============================================================
\begin{abstract}
We present \textbf{Floodcaster}, a mobile operator system that
demonstrates the FloodRSCT certificate-based diagnostic framework
\citep{martin2026floodrsct} in a live emergency-management workflow.
For a 72-hour flood scenario, Floodcaster loads model forecasts,
computes location-level reliability certificates, and maps each
certificate to a Trust, Review, Escalate, or Suppress action.
Supporting evidence is exposed through three operator surfaces:
an action queue, a four-level drill hierarchy, and an evidence
drawer with authority-tier and provenance visibility. Every action
is backed by a snapshot-immutable audit trace, replayable
bit-identically from a single record. The demonstration walks
attendees through one urban and one storm-surge scenario, runs a
live stress-test recipe, and replays an audit trace end-to-end.
\end{abstract}

\keywords{flood crisis decision support, mobile operator systems,
  auditable AI, model risk, certificate-based diagnostics,
  emergency management}

\maketitle

% ============================================================
% 1. INTRODUCTION (≈0.5 page)
% ============================================================
\section{Introduction}
\label{sec:demo-intro}

Flood emergency managers receive risk scores; they need decisions.
A 72-hour flood forecast that ranks several neighborhoods as high risk
does not tell an operator which warnings are reliable, which need
analyst review, and which should be suppressed because the model is
outside its operating zone. \textbf{Floodcaster} is the mobile
operator system that closes this loop. It loads scenario forecasts,
applies the FloodRSCT readiness
certificate~\citep{martin2026floodrsct}, and presents the resulting
action queue in a form an emergency manager can act on in the field.

This demonstration paper focuses on the \emph{operator-facing
artifact}: the mobile surfaces, the drill hierarchy, the audit trace,
and the live workflow. The underlying certificate machinery (signal
quality $\alpha$, representation--solver compatibility $\kappa$,
representational turbulence $\sigma$, task-difficulty bound
$\mathrm{TRF}$, geometry flags, and the fixed-order action gate)
is presented in the companion Applications Track
paper~\citep{martin2026floodrsct}. This paper assumes that machinery
and shows what a reviewer can install, operate, and audit.

Three properties distinguish Floodcaster from existing flood
visualization dashboards. First, the action vocabulary is
\textbf{four-valued} (Trust, Review, Escalate, Suppress) rather than
a single scalar priority score. Second, every action is backed by an
\textbf{auditable evidence chain}: the operator can drill from the
action chip into the certificate fields, into the retrieved evidence,
and into the provenance DAG of each evidence item. Third, the entire
session is \textbf{snapshot-immutable}: re-running any query at the
original snapshot returns bit-identical evidence, which is what makes
after-action review and regulatory replay possible.

% ============================================================
% 2. SYSTEM OVERVIEW (≈0.75 page)
% ============================================================
\section{System Overview}
\label{sec:demo-system}

Floodcaster is structured as a four-layer pipeline
(Figure~\ref{fig:demo-arch}; details
in~\citet{martin2026floodrsct}): a \emph{forecast layer} that emits
per-location risk scores, a \emph{certificate layer} that diagnoses
forecast readiness, a \emph{retrieval layer} that supplies
authority-ranked grounding evidence, and a \emph{crisis-action layer}
that maps the certificate to an operational action via a fixed-order
gate. The mobile application is the operator-facing instantiation of
the crisis-action layer, with downward visibility into all three
upstream layers.

\begin{figure}[h]
\centering
% \includegraphics[width=\columnwidth]{figures/floodcaster_mobile_arch.pdf}
\fbox{\rule{0pt}{3.5cm}\hspace{0.92\columnwidth}}
\caption{Floodcaster mobile system. Top: action queue keyed to the
  current scenario snapshot. Middle: drill view exposing the
  certificate fields, the gate sequence, the retrieved evidence (with
  R/S/N annotations and Tier 1/2/3 badges), and the provenance DAG.
  Bottom: audit-trace footer with a one-tap replay handle.}
\Description{A three-pane mobile screen mockup. The top pane shows
  an action queue listing locations with color-coded
  Trust, Review, Escalate, or Suppress chips and compliance badges. The
  middle pane shows a drilled-in location view exposing certificate
  fields, the gate sequence, retrieved evidence items with authority
  tier badges and R/S/N scores, and a small provenance graph. The
  bottom pane shows an audit-trace footer with an export-and-replay
  button.}
\label{fig:demo-arch}
\end{figure}

\paragraph{Backend.}
Floodcaster is backed by a Postgres evidence store with the pgvector
extension. Each evidence object is keyed by a composite snapshot
identifier (snapshot timestamp, embedding model version, index
version) that enforces bit-identical replay. Authority tier (Tier 1
regulatory, Tier 2 operational SOP, Tier 3 technical), source family,
and gate-question-relative $(R, S_{\text{sup}}, N)$ annotations are
first-class columns rather than post-hoc metadata.

\paragraph{Scenario routing.}
The crisis-action layer consults a deterministic scenario router that
declares, per scenario and per gate question, which evidence families
are \emph{allowed}, \emph{required}, and \emph{forbidden}. No language
model sits in the routing path; routing decisions are inspectable and
replayable, consistent with model-risk
expectations~\citep{sr11-7,nistairmf}.

\paragraph{Five flood scenarios on demo day.}
The mobile artifact ships pre-loaded with five scenarios---Houston,
New Orleans, Riverside-Coachella, Southwest Florida, and NYC/NJ---each
with a representative 72-hour snapshot and its own hydrometeorological
character (urban pluvial, protected-basin infrastructure, desert flash
flood, storm surge, and dense urban cloudburst respectively; see
\citet{martin2026floodrsct} for the per-scenario data registries).
Attendees can swap scenarios from the home screen and watch the
action queue re-populate from the new snapshot.

% ============================================================
% 3. MOBILE OPERATOR SURFACES (≈1 page)
% ============================================================
\section{Mobile Operator Surfaces}
\label{sec:demo-surfaces}

The mobile interface exposes the FloodRSCT pipeline through three
primary surfaces and three drill-down panels.

\paragraph{Action Queue (home screen).}
The home screen lists every location in the current scenario sorted
by action priority. Each row carries the location identifier, the
risk score, the assigned action with an associated color chip
(Trust, Review, Escalate, or Suppress), and a compliance badge
indicating whether the grounding gate was supported by Tier-1
evidence under the scenario's registry. Locations in Suppress do
\emph{not} show a risk numeral as the primary affordance; the action
is the primary affordance, the score is secondary. This is a
deliberate design choice that prevents an operator from acting on a
high-risk number that the system itself flagged as unreliable.

\paragraph{Drill Hierarchy.}
Tapping a location opens the drill view, which exposes a four-level
hierarchy: county $\to$ ZCTA $\to$ modal source $\to$ feature. Modal
sources are scenario-specific:
drainage links and stream gauges in Houston, levee segments and pump
stations in New Orleans, canyon routing and road-access features in
Riverside--Coachella, surge zones and evacuation routes in Southwest
Florida, sewer sheds and subway stations in NYC/NJ. The drill view
is implemented as a single recursive component; the scenario
registry determines what is shown at each level.

\paragraph{Evidence Drawer.}
A drawer attached to the drill view shows the ranked evidence items
retrieved for the location's gate questions, each item annotated with
its authority tier (Tier~1 / Tier~2 / Tier~3), its
gate-question-relative $(R, S_{\text{sup}}, N)$ scores, and a one-tap
expansion to the source document. Evidence items that satisfy a
\emph{forbidden} registry slot (proxy-only justification for a Gate~4
firing, for example) are shown but visually marked as
non-substantiating, with an explanatory annotation. This is what
makes the grounding contract visible at the operator level rather
than hidden in the certificate fields.

\paragraph{Provenance Panel.}
Each evidence item exposes a small DAG (not a single URL) showing the
transform chain from raw source to the variable used in the
certificate---for example, raw NWS QPF tiles $\rightarrow$ basin
aggregation $\rightarrow$ ZCTA crosswalk $\rightarrow$ rainfall
intensity feature. After-action auditors need this chain, not a
hyperlink to a press release.

\paragraph{Compliance Badge.}
Attached to each action chip, the compliance badge is a one-glance
indicator of whether the grounding gate fired under Tier-1 precedence
required by the scenario's registry. Three states: green (Tier-1
supported), amber (Tier-2/3 only), red (forbidden-family
substitution attempted, gate failed closed). The badge is the
shortest path between an operator and a defensible decision.

\paragraph{Audit-Trace Footer.}
Every action carries a footer with a one-tap audit handle. The handle
yields a structured record (\S\ref{sec:demo-audit}) that uniquely
identifies the certificate, the evidence bundle, and the gate
sequence that produced this action at this snapshot.

% ============================================================
% 4. AUDIT TRACE AND REPLAYABILITY (≈0.5 page)
% ============================================================
\section{Audit Trace and Replayability}
\label{sec:demo-audit}

The Floodcaster audit record is a single JSON object with the
following fields:

\begin{itemize}
\setlength{\itemsep}{1pt}
\item \texttt{scenario\_id}, \texttt{location\_id}, \texttt{snapshot\_t}
  --- compose the snapshot-immutability key.
\item \texttt{model\_id}, \texttt{representation\_id} --- identify
  the forecast layer.
\item \texttt{risk\_score} --- the raw model output.
\item \texttt{certificate\_fields} --- the
  $\alpha, \kappa, \sigma, \mathrm{TRF}$ tuple together with
  geometry flags.
\item \texttt{gate\_sequence} --- the ordered firings through the
  fixed-order action gate.
\item \texttt{evidence\_bundle} --- an array of retrieved evidence
  object IDs. Each item carries gate-question-relative
  $(R, S_{\text{sup}}, N)$ scores and a Tier~1/2/3 authority badge.
\item \texttt{action}, \texttt{rationale\_text} --- the final
  assigned action and its generated explanation.
\item \texttt{tier1\_compliance} --- binary indicator for whether
  Gate~4 was Tier-1 supported.
\item \texttt{timestamp} --- record creation time.
\end{itemize}

This is sufficient to replay any action bit-identically. The
demonstration includes a live replay: an audit record exported from
the mobile device is opened on the demo laptop, the snapshot key is
extracted, the same evidence bundle is retrieved from the Postgres
backend, the certificate is recomputed, and the action chip is
regenerated. Round-trip integrity is a property of the snapshot key,
not a guarantee we can offer here in prose; the demo shows it.

% ============================================================
% 5. LIVE DEMONSTRATION PLAN (≈0.75 page)
% ============================================================
\section{Live Demonstration Plan}
\label{sec:demo-plan}

The demonstration is designed for a 12--15 minute slot per attendee
group, with the following script:

\paragraph{Act 1 --- Urban pluvial scenario (Houston).}
Open the Houston/Harris County snapshot. Show the action queue with
a mix of Trust, Review, Escalate, or Suppress assignments. Drill into one
Escalate location, walk the evidence drawer (illustrating Tier-1
NFHL + Tier-3 USGS gauge support), and show the provenance DAG for
the rainfall feature. Highlight one Suppress location and explain
why the certificate flagged it (typically: proxy-driven Gate~4
firing detected, forbidden-family substitution attempted).

\paragraph{Act 2 --- Storm surge scenario (Southwest Florida).}
Swap to the SW Florida snapshot. Demonstrate scenario-specific
modal sources (surge zones, evacuation routes) and show how the
\emph{same} mobile component, fed the \emph{same} certificate, with
the \emph{same} gate order, produces a scenario-appropriate action
distribution. This is the visible payoff of the single-compute-path
design: one app, five scenarios, no special-casing.

\paragraph{Act 3 --- Live stress test.}
From the drill view, invoke a stress-test recipe (sensor dropout,
rainfall shift, or infrastructure failure). The action queue
re-populates with the perturbed actions; the compliance badges shift
visibly on several locations. Discuss the
Trust $\to$ Review, Review $\to$ Escalate, and
Escalate $\to$ Suppress drift patterns.

\paragraph{Act 4 --- Audit replay.}
Export the audit record for one of the perturbed actions. Open it on
the demo laptop. Run the replay command. Show that the same evidence
bundle and the same action are reproduced bit-identically. Close
with a one-slide summary of the four-layer separation and the role
each layer plays in the audit trail.

\paragraph{Requirements at the venue.}
A standard demo table with power, a small monitor (mirroring the
mobile screen for group viewing) connected via HDMI or USB-C, and
wired or stable wireless networking for the backend connection. The
five demo snapshots are cached locally; no external network is
required for the certificate or replay components.

% ============================================================
% 6. RELATIONSHIP TO PRIOR DEMOS AND CONCLUSION
% ============================================================
\section{Relation to Prior Demos and Conclusion}
\label{sec:demo-related}

Prior SIGSPATIAL demonstrations of flood and emergency management
systems have emphasized real-time data ingestion, visualization
density, and routing optimization. Floodcaster is complementary: it
takes flood model outputs as input and demonstrates how to
\emph{interpret} them defensibly. The contribution of this
demonstration is an operator-facing instantiation of a
certificate-grounded crisis decision workflow, with full audit
replay, on a mobile form factor that an emergency manager can carry
into the field. We invite attendees to operate the system, attempt
to break it with stress-test recipes, and inspect the audit trace
from any action they generate.

% ============================================================
% REFERENCES
% ============================================================
\bibliographystyle{ACM-Reference-Format}
\bibliography{floodcaster_demo}

\end{document}
